% Proposed manuscript insertion. See integration/README.md before installing.
% Uses the manuscript's existing theorem environments; no global macros added.
\section{Integral distribution normal forms and reflection torsion}
\label{intdist:sec:integral}
The characteristic-zero conductor decomposition can be strengthened to an
integral original-symbol statement. Reflection then introduces a separate,
exactly computable two-torsion phenomenon. We record the results of the
integral continuation~\cite{IntegralDistributionReflection2026}; that report
contains complete proofs, an explicit resolution, and independent certificates.

For $q\ge2$, let $P$ be its prime divisors and put
\[
 R=\mathbb Z[A_p:p\in P],\qquad X_q=(q^{-1}\mathbb Z)/\mathbb Z.
\]
Let $\mathcal U_q$ be the free $R$-module on $e_x$, $x\in X_q$, modulo
all rows
\[
 r_{p,x}=\sum_{py=x}e_y-A_pe_x,\qquad p\mid q,\quad x\in X_{q/p}.
\]
Composite weights are products of the prime weights. The endpoint $e_0$
is a formal symbol; under polylogarithm evaluation it means $\zeta(s)$.

For $q=\prod p^{e_p}$, write uniquely
$x=\sum a_p/p^{e_p}$ with $0\le a_p<p^{e_p}$. Retain the point $x$ when
\[
 a_p\notin\{1,\ldots,p^{e_p-1}\}\quad\hbox{for every }p\mid q,
\]
and denote the resulting set by $\mathcal B_q$. It has $\varphi(q)$
points and contains zero.

\begin{theorem}[Integral original-symbol basis]
\label{intdist:thm:unit}
The classes $e_b$, $b\in\mathcal B_q$, form a basis of $\mathcal U_q$
over $R$. A fixed subset of $q-\varphi(q)$ original prime rows has a
unit lower triangular submatrix on the complementary point columns.
The entire raw matrix is equivalent over $R$ to an identity block of
size $q-\varphi(q)$ and zeros. These assertions survive arbitrary base
change to every commutative coefficient ring.
\end{theorem}
\begin{proof}[Proof outline]
For a nonretained point choose its least forbidden prime $p$ and the row
$r_{p,px}$. Order points by total exact prime-power height, then by the
number of forbidden coordinates. Every other point in this row precedes
$x$, and the coefficient of $e_x$ is one. This gives the triangular unit
minor and a free quotient by the selected rows. The complete conductor
calculation over the fraction field gives rank $q-\varphi(q)$ for all
prime rows. The selected quotient is torsion-free over $R$, so every
remaining raw row already vanishes integrally in that quotient. The
resulting polynomial matrix equivalence survives every base change.
\end{proof}

The associated normal forms are computed without division:
\[
 \operatorname{NF}(x)=A_p\operatorname{NF}(px)
       -\sum_{j=1}^{p-1}\operatorname{NF}(x+j/p)
\]
for a nonretained point, and $\operatorname{NF}(b)=e_b$ on the basis.
This determines explicit integer polynomials $P_{x,b}(\boldsymbol A)$.
It proves that the unreflected integer Smith factors are all one,
including expanded finite-jet matrices. A character Fourier transform
is not needed in the actual certificate verifier.

\subsection{Reflection is a different integral quotient}
Let $J e_x=e_{-x}$ and define
\[
 C_+(U)=U/(J-1)U,\qquad C_-(U)=U/(J+1)U.
\]
These are coinvariants, not integral eigensublattices. Put
\[
 r_* = \#\{p\mid q:p\text{ odd}\}+\boldsymbol1_{4\mid q}.
\]
\begin{theorem}[Scalar reflected Smith factors]
\label{intdist:thm:scalar}
For $q>2$ and arbitrary integer weights $A_p=a_p$,
\[
 C_\pm\simeq\mathbb Z^{\varphi(q)/2}\oplus(\mathbb Z/2\mathbb Z)^h,
 \qquad
 h=\begin{cases}
 0,&\text{some odd-prime weight is even},\\
 2^{r_*-1},&\text{all odd-prime weights are odd}.
 \end{cases}
\]
For $q=2$, $C_+=\mathbb Z$ and $C_-=\mathbb Z/2\mathbb Z$.
\end{theorem}

A full proof uses the weighted Anderson resolution with terms
$\bigoplus_{|S|=k}R[X_{q/\prod_{p\in S}p}]$.
The two-periodic reflection complex has an explicit acyclic subcomplex
consisting of nonfixed point coordinates and even multiples of fixed
coordinates in even degree, and all coordinates in odd degree.
The remaining fixed-point complex is the binary Koszul complex on
\[
 1-\overline a_p\quad(p\text{ odd}),\qquad
 \overline a_2(1-\overline a_2)\quad\text{when }4\mid q.
\]
Its odd homology gives the torsion of $C_+$ and its even homology gives
that of $C_-$. This also explains why a lone factor of two in the level
does not increase $r_*$.

\begin{theorem}[One-variable integral jets]
\label{intdist:thm:jets}
Let $q>2$, $B_L=\mathbb Z[u]/(u^L)$, and take arbitrary weights
$a_p(u)\in B_L$. Reduce the preceding Koszul parameters modulo two,
and let $v$ be their minimum truncated $u$-valuation, with the zero
polynomial assigned valuation $L$. Then, as abelian groups,
\[
 C_\pm\simeq\mathbb Z^{L\varphi(q)/2}
       \oplus(\mathbb Z/2\mathbb Z)^{v2^{r_*-1}}.
\]
Their integer-torsion submodules, as $B_L$-modules, are each
$(B_L/(2,u^v))^{2^{r_*-1}}$.
\end{theorem}
Indeed, over $\mathbb F_2[u]/(u^L)$ the parameter vector is equivalent
to $(u^v,0,\ldots,0)$. Its Koszul homology has
$\binom{r_*}{k}$ copies of $\mathbb F_2[u]/(u^v)$ in degree $k$.
The free ranks follow from the zero rational reflection trace.
The free abelian part need not be a free $B_L$-module; no such claim
is part of the theorem.

\subsection{An exact all-order identity}
Substitute $A_p=p^{1-s}$ into the polynomial normal forms. Root filtering
and meromorphic continuation prove their evaluation for every complex
order. A three-row certificate at level twelve yields
\begin{align}
 &\Li_s(e^{\pi i/6})+\Li_s(e^{5\pi i/6})
       +(1+3^{1-s})\Li_s(-i)\nonumber\\
 &\hspace{35mm}=(12^{1-s}-6^{1-s})\zeta(s).
 \label{intdist:eq:level12}
\end{align}
The apparent pole at $s=1$ is removable, with value $-\log2$.
The continuation gives all order derivatives there by Laurent
multiplication; the first derivative of the right side is
\[
 \frac{\log^2 12-\log^2 6}{2}-\gamma\log2.
\]
In general, differentiating an endpoint coefficient that vanishes at
$s=1$ produces an extra contribution from its derivative of order
one higher. The report states the complete finite-part formula and an
analogous five-row identity at level thirty.

\paragraph{Scope of the completion.}
These results settle the integral basis, unreflected Smith, scalar
reflection, and one-variable integral-jet questions for this precise
presentation. Multivariate reflection is reduced to Koszul homology,
not to an unjustified single valuation. Arithmetic independence of the
evaluated periods, higher-depth relations, and the Gaussian $S_6$
conjecture remain distinct questions. The ordinary distribution and
Anderson resolution have classical antecedents; the report preserves
the relevant attributions and does not claim worldwide priority for
all of the displayed consequences.
